\documentclass[11pt]{article}

\usepackage[margin=1in]{geometry}
\usepackage{amsmath,amssymb,amsthm,amsfonts}
\usepackage{mathtools}
\usepackage{bm}
\usepackage{hyperref}
\usepackage{enumitem}
\usepackage{graphicx}
\usepackage{physics}
\usepackage{bbm}

\title{Multiplicity Prime Dynamics:\\
Hyperprime Tensor Evolution and Prime-Attention Neural Layers}
\author{Citizen Gardens \\ \\ The Foundation of Multiplicity}
\date{April 5, 2026}

\begin{document}
\maketitle

\begin{abstract}
This report synthesizes and formalizes a line of developments around Multiplicity Prime Dynamics (MPD), starting from a hyperprime tensor framework and culminating in concrete, executable models: a finite $\{2,3\}$ tensor network and a prime-attention neural layer with an associated PyTorch experiment. The core novelty is a closed feedback loop in which prime-indexed operators simultaneously generate dynamics, structure Hilbert evolution, govern cognitive metrics, and define ontological invariants. We provide an executive summary, a comprehensive mathematical overview, explicit layer and dynamics equations, and an experiment specification suitable for immediate numerical validation.
\end{abstract}

\tableofcontents

\section{Executive Summary}

This document consolidates a progression of ideas that began with a hyperprime tensor evolution framework and moved toward a fully testable research program under the umbrella of \emph{Multiplicity Prime Dynamics} (MPD). The key developments are:

\begin{itemize}[nosep]
  \item A clarification of what is truly novel: not the individual ingredients (tensor networks, $p$-adic norms, gauge fields, categorical structure), but the way \emph{primes} act as a unifying operator family across dynamics, cognition, and ontology.
  \item A structural repair of earlier conceptual tensions: primes become spectral constraints rather than mutable fields; their roles are strictly stratified (generator, operator, spectrum, ontology); the $p$-adic metric is realized as a finite prime-fiber bundle; and ``ontology'' is reinterpreted as membership in invariant subspaces.
  \item A concrete finite model: a 4-qubit tensor network with active primes $\{2,3\}$, noncommuting prime operators $\hat{O}_2,\hat{O}_3$, a prime spectrum operator $\hat{P}$, and time evolution $\Xi(t)$ with prime-dependent attention weights.
  \item A prime-attention neural layer: a finite-dimensional, trainable component where branch weights $W_p$, ontology projectors $\Pi_p$, prime attention $w_p$, and fiber coordinates $c_p$ instantiate the MPD feedback loop inside a standard ML setting.
  \item A ready-to-run PyTorch experiment: synthetic prime-colored data, training of the prime-attention layer, logging of prime usage and fiber norms, a cognitive distance $d_C$ for clustering, and a path-dependence test via representation divergence across prime orderings.
\end{itemize}

The result is a coherent theory with explicit equations and immediate numerical tests, transforming a philosophically rich vision into a falsifiable computational research program.

\section{Conceptual and Structural Overview}

\subsection{Core MPD Idea}

Multiplicity Prime Dynamics (MPD) posits that \emph{primes} should be treated as an integrated operator family:

\begin{itemize}[nosep]
  \item They \emph{generate dynamics} through prime-indexed operators $\hat{O}_p^{(r)}$ acting on multiplicity spaces.
  \item They \emph{parameterize Hilbert evolution} via prime-decomposed flows $\Xi(t)$.
  \item They \emph{structure cognition} through prime-resolved metrics and fiber bundles.
  \item They \emph{define ontology} as invariant subspaces of prime dynamics.
\end{itemize}

These roles are glued into a single feedback loop:
\[
\text{prime spectrum} \;\leftrightarrow\; \text{tensor evolution} \;\leftrightarrow\; \text{cognitive metric} \;\leftrightarrow\; \text{ontology} \;\leftrightarrow\; \text{prime spectrum}.
\]

The novelty lies in this closure, not in any single ingredient.

\subsection{Role Hierarchy for Primes}

To avoid semantic overload, MPD enforces a strict hierarchy:

\begin{center}
\begin{tabular}{ll}
\textbf{Role} & \textbf{Meaning} \\
\hline
prime & label for an operator family ($p \in \mathbb{P}$) \\
operator & dynamical action $\hat{O}_p$ on a space \\
spectrum & eigenvalues of $\hat{P}$, constrained by primes \\
ontology & invariant structures under $\hat{O}_p$ (fixed-point subspaces)
\end{tabular}
\end{center}

Primes themselves are no longer treated as evolving field values but as spectral constraints encoded in an operator
\[
\hat{P} = \sum_p p\,\Pi_p,
\]
where $\Pi_p$ are spectral projectors.

\section{Mathematical Framework}

\subsection{Multiplicity Spaces and Prime Operators}

Let $\mathcal{H}$ be a (finite-dimensional) Hilbert space equipped with a family of \emph{prime operators}
\[
\hat{O}_p \in \mathrm{End}(\mathcal{H}),\qquad p \in P \subset \mathbb{P}.
\]

In the finite model, $P$ is a small active prime set, e.g.\ $P = \{ 2,3\}$ or $P = \{2,3,5\}$. Each $\hat{O}_p$ is typically Hermitian (in the tensor model) or a real linear map (in the neural layer).

\subsection{Prime Spectrum Operator}

Each $\hat{O}_p$ admits an eigendecomposition
\[
\hat{O}_p = \sum_{\lambda} \lambda \,\Pi_{p,\lambda},
\]
with $\Pi_{p,\lambda}$ the projectors onto eigen-subspaces. We select a ``dominant'' eigenvalue $\lambda_{\max}^{(p)}$ and define
\[
\Pi_p := \Pi_{p,\lambda_{\max}^{(p)}},
\]
and the prime spectrum operator
\[
\hat{P} = \sum_{p\in P} p\, \Pi_p.
\]

Here, primes appear as weights on spectral projectors. Dynamics does not change $p$; instead, it acts within the spaces on which $\Pi_p$ project. The noncommutativity
\[
[\hat{P}, \hat{O}_p] \neq 0
\]
encodes the idea that prime structure is not trivial with respect to the dynamics.

\subsection{Prime-Decomposed Evolution}

In MPD, evolution is expressed as a prime-decomposed flow
\[
\Xi(t) = \sum_{p\in P} w_p(t)\,U_p(t),
\]
where $U_p(t)$ are prime-specific evolutions and $w_p(t)$ are attention-like weights.

In a discrete-time setting with small timestep $\Delta t$, one can take
\[
U_p(t) = \exp\!\left(-i \hat{O}_p\Delta t\right)
\]
in the quantum/tensor context, or a general parametrized linear map in the neural context.

The state $x_t \in \mathcal{H}$ evolves as
\[
x_{t+1} = \frac{\Xi(t) x_t}{\|\Xi(t) x_t\|}.
\]

\subsection{Cognitive Bundle and Metric}

Cognitive states are modeled as prime-fiber coordinates over the active prime set $P$. For a given state $x \in \mathcal{H}$, define
\[
c_p(x) := \Pi_p x \in \mathcal{H}, \qquad p \in P,
\]
and the bundle
\[
\mathcal{H}_{\mathrm{cog}} = \prod_{p\in P} \mathcal{H}.
\]

The cognitive distance between two states $x,y$ is then modeled as a sup-norm over the fibers:
\[
d_C(x,y) = \max_{p\in P} \|c_p(x) - c_p(y)\|_2.
\]

This is a finite-prime analogue of a $p$-adic sup metric, made explicitly computable.

\subsection{Ontology via Fixed-Point Projectors}

Ontological questions about ``whether prime $p$ is cosmic'' for a state $x$ are recast as fixed-point constraints:

\[
\hat{Q}_p(x) := \mathrm{Proj}_{\mathrm{Fix}(\hat{O}_p)}(x),
\]
where $\mathrm{Fix}(\hat{O}_p)$ is the invariant subspace of $\hat{O}_p$. The associated observable
\[
q_p(x) := \langle x,\hat{Q}_p(x)\rangle
\]
measures the degree to which $x$ lies in the $p$-invariant regime.

\section{Finite $\{2,3\}$ Tensor Network Model}

\subsection{Hilbert Space and Pauli Lifting}

Consider a 4-qubit system
\[
\mathcal{H} = (\mathbb{C}^2)^{\otimes 4},\quad \dim(\mathcal{H})=16.
\]

Let $X,Y,Z$ be Pauli matrices on a single qubit. For $k=1,\ldots,4$, denote by $X_k,Y_k,Z_k$ the corresponding operators acting on qubit $k$ (via tensoring with identity on the others).

\subsection{Prime Operators with Overlap}

Define noncommuting Hermitian operators $\hat{O}_2,\hat{O}_3$ acting on overlapping qubits:
\begin{align*}
\hat{O}_2 &= X_1 Z_2 + Z_1 X_2 + \gamma\, Z_1 Z_2,\\
\hat{O}_3 &= Y_2 X_3 + X_2 Y_3 + \gamma\, Z_2 Z_3,
\end{align*}
with coupling parameter $\gamma>0$. These share qubit 2, ensuring that generically
\[
[\hat{O}_2,\hat{O}_3] \neq 0.
\]

\subsection{Spectral Projectors and Prime Spectrum}

Diagonalize each operator:
\[
\hat{O}_2 = \sum_\lambda \lambda\, \Pi_{2,\lambda}, \quad
\hat{O}_3 = \sum_\mu \mu\, \Pi_{3,\mu}.
\]

Select largest-magnitude eigenvalues $\lambda_{\max},\mu_{\max}$ and define
\[
\Pi_2 = \Pi_{2,\lambda_{\max}},\quad
\Pi_3 = \Pi_{3,\mu_{\max}}.
\]

Prime spectrum operator:
\[
\hat{P} = 2\,\Pi_2 + 3\,\Pi_3.
\]

\subsection{Time Evolution and Prime Attention}

Let
\[
U_2 = \exp(-i \hat{O}_2 \Delta t),\quad
U_3 = \exp(-i \hat{O}_3 \Delta t)
\]
for a fixed $\Delta t>0$.

Given a state vector $|\psi(t)\rangle \in \mathcal{H}$, define prime overlaps
\[
a_p(t) = \|\Pi_p |\psi(t)\rangle\|^2 = \langle \psi(t)|\Pi_p|\psi(t)\rangle,\quad p\in\{2,3\}.
\]

Define attention weights via softmax with inverse temperature $\beta>0$:
\[
w_p(t) = \frac{\exp(\beta a_p(t))}{\exp(\beta a_2(t))+\exp(\beta a_3(t))}.
\]

Construct the effective evolution operator
\[
\Xi(t) = w_2(t) U_2 + w_3(t) U_3,
\]
and evolve
\[
|\psi(t+1)\rangle = \frac{\Xi(t)|\psi(t)\rangle}{\|\Xi(t)|\psi(t)\rangle\|}.
\]

This defines a nonlinear, prime-attended discrete dynamics.

\subsection{Observables and Predictions}

Key observables:

\begin{itemize}[nosep]
  \item Ontology expectations: $q_p(t) = \langle \psi(t)|\Pi_p|\psi(t)\rangle$.
  \item Prime attention weights: $w_p(t)$.
  \item Memory of initial state: fidelity $F_0(t) = |\langle \psi(0)|\psi(t)\rangle|^2$.
  \item Spectral statistics of $\Xi(t)$: eigenvalues in the complex plane.
\end{itemize}

Predictions:

\begin{itemize}[nosep]
  \item \textbf{Spectral clustering:} eigenvalues of $\Xi(t)$ cluster near the spectra of $U_2$ and $U_3$, modulated by $w_p(t)$.
  \item \textbf{Log-periodicity:} oscillatory patterns in $q_p(t)$ as a function of $\log t$, suggestive of discrete scale invariance.
  \item \textbf{Phase transitions:} as $\gamma$ or $\beta$ vary, the system may transition between regimes dominated by prime 2 or 3.
  \item \textbf{Path dependence:} differing sequences of $U_2,U_3$ applications result in divergent final states when $[\hat{O}_2,\hat{O}_3]\neq 0$.
\end{itemize}

\section{Prime-Attention Neural Layer}

\subsection{Layer Structure}

Let $x \in \mathbb{R}^{d_{\mathrm{in}}}$ be an input and $h \in \mathbb{R}^{d_{\mathrm{out}}}$ an output. Fix an active prime set $P=\{2,3,5\}$.

For each $p\in P$:

\begin{itemize}[nosep]
  \item A weight matrix $W_p \in \mathbb{R}^{d_{\mathrm{out}}\times d_{\mathrm{in}}}$.
  \item A learnable basis $U_p \in \mathbb{R}^{d_{\mathrm{out}}\times k_p}$, with columns orthonormalized to form
  \[
  \Pi_p = U_p U_p^\top.
  \]
\end{itemize}

Let $\sigma$ be a nonlinearity (e.g.\ ReLU).

\subsection{Branch Responses and Ontology Overlaps}

Prime branch responses:
\[
z_p = W_p x + b_p,\qquad \tilde{z}_p = \sigma(z_p).
\]

Prime ontology projections:
\[
c_p(x) = \Pi_p \tilde{z}_p,\quad a_p(x) = \|c_p(x)\|_2^2 = \tilde{z}_p^\top \Pi_p \tilde{z}_p.
\]

\subsection{Prime Attention and Output}

Attention weights:
\[
w_p(x) = \frac{\exp(\beta a_p(x))}{\sum_{q\in P} \exp(\beta a_q(x))}.
\]

Layer output:
\[
h(x) = \sum_{p\in P} w_p(x)\,\tilde{z}_p.
\]

Optionally, a final nonlinearity $\phi$ can be applied:
\[
y(x)=\phi(h(x)).
\]

\subsection{Cognitive Bundle Interpretation}

The collection of projected fibers
\[
C(x) = \big(c_p(x)\big)_{p\in P} \in \prod_{p\in P}\mathbb{R}^{d_{\mathrm{out}}}
\]
realizes a finite-dimensional cognitive bundle. A cognitive distance between activations $x,x'$ is defined as
\[
d_C(x,x') = \max_{p\in P} \|c_p(x) - c_p(x')\|_2.
\]

\subsection{Training and Regularization}

Consider a classifier network:
\[
x \mapsto h_1(x) \mapsto h_2(x) \mapsto \hat{y},
\]
where $h_1,h_2$ are prime-attention layers, and
\[
\hat{y} = \mathrm{softmax}(W_{\mathrm{cls}} h_2 + b_{\mathrm{cls}}).
\]

For labeled data $(x_i,y_i)$, the classification loss is
\[
\mathcal{L}_{\mathrm{cls}} = -\sum_i \sum_k \mathbbm{1}_{\{y_i=k\}} \log \hat{y}_{i,k}.
\]

To encourage MPD-like behavior, we add:

\paragraph{Sparsity regularizer.} Encourage near one-hot prime attention:
\[
\mathcal{L}_{\mathrm{sparse}} = \lambda_s\,\mathbb{E}_x\left[\sum_{p\in P} w_p(x)\big(1-w_p(x)\big)\right],
\]
minimized when attention is concentrated on a single prime.

\paragraph{Alignment regularizer.} Encourage branch outputs to align with their own subspaces:
\[
\mathcal{L}_{\mathrm{align}} = \lambda_a\,\mathbb{E}_x\left[\frac{1}{|P|}\sum_{p\in P}\|\tilde{z}_p - c_p(x)\|_2^2\right].
\]

Total loss:
\[
\mathcal{L} = \mathcal{L}_{\mathrm{cls}} + \mathcal{L}_{\mathrm{sparse}} + \mathcal{L}_{\mathrm{align}}.
\]

\section{PyTorch Implementation Snippet}

Below is a schematic version of the prime-attention layer in PyTorch form, reflecting the equations above (omitting imports and training loop for brevity):

\begin{verbatim}
class PrimeAttentionLayer(nn.Module):
    def __init__(self, d_in, d_out, primes=(2,3,5), k_sub=8, beta=5.0):
        super().__init__()
        self.primes = primes
        self.beta = beta
        self.d_out = d_out

        self.W = nn.ModuleDict({
            str(p): nn.Linear(d_in, d_out, bias=True) for p in primes
        })
        self.U = nn.ParameterDict({
            str(p): nn.Parameter(torch.randn(d_out, k_sub)) for p in primes
        })

    def forward(self, x, return_fibers=False):
        branch_outputs = {}
        fibers = {}
        overlaps = []

        for p in self.primes:
            z = self.W[str(p)](x)
            z = F.relu(z)
            branch_outputs[p] = z

            U = self.U[str(p)]
            Q, _ = torch.linalg.qr(U)
            Pi = Q @ Q.T

            c = z @ Pi.T
            fibers[p] = c
            a = (c * c).sum(dim=1)
            overlaps.append(a.unsqueeze(1))

        overlaps_cat = torch.cat(overlaps, dim=1)
        weights = F.softmax(self.beta * overlaps_cat, dim=1)

        h = torch.zeros(x.size(0), self.d_out, device=x.device)
        for i, p in enumerate(self.primes):
            w_p = weights[:, i].unsqueeze(1)
            h = h + w_p * branch_outputs[p]

        if return_fibers:
            return h, weights, fibers
        return h, weights
\end{verbatim}

A full script can include:

\begin{itemize}[nosep]
  \item Synthetic prime-colored data generation (each class produced by a different prime-linear map).
  \item Training loop with logging of $w_p$ and $\|c_p\|^2$ per epoch.
  \item Cognitive distance clustering using $d_C$ and comparison to Euclidean clustering.
  \item Path dependence test by training two-layer networks with differing prime compositions per layer and comparing hidden representations via CKA.
\end{itemize}

\section{Recommended Experimental Protocol}

To interpret numerical experiments in MPD terms:

\begin{enumerate}[nosep]
  \item \textbf{Prime dominance trajectories.} Plot the epoch-wise averages $\bar{w}_p = \mathbb{E}_x[w_p(x)]$ to detect phase transitions from mixed prime usage to prime dominance as hyperparameters such as $\beta$ vary.
  \item \textbf{Prime specialization.} Track $\mathbb{E}_x\|c_p(x)\|^2$ per class to see which primes specialize to which data regimes.
  \item \textbf{Cognitive clustering.} Compare cluster quality (e.g.\ Adjusted Rand Index) under Euclidean distance in hidden space vs.\ $d_C$ in fiber space.
  \item \textbf{Path dependence.} Train networks with different prime-layer compositions and compare representation similarity (e.g.\ via CKA) to show noncommutative path dependence even under similar classification accuracy.
\end{enumerate}

These steps connect empirical signatures directly to the theoretical claims of MPD.

\section{Conclusion}

We have moved from a high-level hyperprime tensor framework to a mathematically disciplined and empirically testable Multiplicity Prime Dynamics program. Primes are now systematically embedded as operator labels, spectral constraints, cognitive fibers, and ontological invariants. Both the finite $\{2,3\}$ tensor network and the prime-attention neural layer provide tractable laboratories where MPD's qualitative claims---phase transitions, discrete scale invariance, prime specialization, and noncommutative path dependence---can be examined quantitatively.

Future work can extend this framework to:

\begin{itemize}[nosep]
  \item richer prime sets and prime tower categories,
  \item continuous-time flows and renormalization-like learning dynamics,
  \item and links to physical tensor network models and $p$-adic/AdS-CFT toy examples.
\end{itemize}

\appendix
\section{Mathematical Appendix: Operator Norm Bounds and Basic Proofs}

In this appendix we collect several elementary but useful estimates and proofs related to the finite $\{2,3\}$ tensor model and the prime-attention neural layer. We work in finite-dimensional Hilbert spaces (or Euclidean spaces), so operator norms are the usual induced norms:
\[
\|A\| := \sup_{\|x\|=1} \|Ax\|.
\]

\subsection{Basic bounds for the finite $\{2,3\}$ tensor model}

Recall the 4-qubit Hilbert space
\[
\mathcal{H} = (\mathbb{C}^2)^{\otimes 4},\quad \dim(\mathcal{H}) = 16,
\]
with prime operators
\begin{align*}
\hat{O}_2 &= X_1 Z_2 + Z_1 X_2 + \gamma\, Z_1 Z_2,\\
\hat{O}_3 &= Y_2 X_3 + X_2 Y_3 + \gamma\, Z_2 Z_3,
\end{align*}
where $X_k,Y_k,Z_k$ are Pauli matrices acting on qubit $k$, and $\gamma>0$ is a coupling parameter.

\begin{lemma}[Norm bounds for prime operators]
\label{lem:OpNormPrime}
With the above notation and the usual operator norm induced by the Hilbert norm on $\mathcal{H}$, we have
\[
\|\hat{O}_2\| \le 2 + |\gamma|,\qquad
\|\hat{O}_3\| \le 2 + |\gamma|.
\]
\end{lemma}

\begin{proof}
Each Pauli matrix $P\in\{X,Y,Z\}$ is unitary with $\|P\|=1$, and tensoring with identities preserves the operator norm. Thus
\[
\|X_1 Z_2\| = 1,\quad \|Z_1 X_2\|=1,\quad \|Z_1 Z_2\|=1,
\]
and similarly for the terms in $\hat{O}_3$. Using the triangle inequality and submultiplicativity of the operator norm, we obtain
\[
\|\hat{O}_2\|
= \|X_1 Z_2 + Z_1 X_2 + \gamma Z_1 Z_2\|
\le \|X_1 Z_2\| + \|Z_1 X_2\| + |\gamma| \,\|Z_1 Z_2\|
\le 1 + 1 + |\gamma| = 2 + |\gamma|.
\]
The bound for $\hat{O}_3$ is identical in form.
\end{proof}

\begin{lemma}[Norm bounds for unitaries]
\label{lem:UnitaryBound}
Let $U_p = \exp(-i \hat{O}_p \Delta t)$ with $\hat{O}_p$ Hermitian. Then $U_p$ is unitary and
\[
\|U_p\| = 1.
\]
\end{lemma}

\begin{proof}
Since $\hat{O}_p^\dagger = \hat{O}_p$, the exponential $U_p = e^{-i \hat{O}_p \Delta t}$ is unitary, hence $U_p^\dagger U_p = I$. For any unitary operator $U$ on a Hilbert space we have $\|U\|=1$, as $\|Ux\|=\|x\|$ for all $x$.
\end{proof}

\begin{lemma}[Norm bound for prime spectrum operator]
\label{lem:PrimeSpectrumBound}
Let $\Pi_2,\Pi_3$ be orthogonal projectors (self-adjoint idempotents) and define
\[
\hat{P} = 2\,\Pi_2 + 3\,\Pi_3.
\]
Then
\[
\|\hat{P}\| \le 2 + 3 = 5.
\]
Moreover, if $\Pi_2$ and $\Pi_3$ have overlapping ranges, then $\|\hat{P}\|\ge 3$.
\end{lemma}

\begin{proof}
Using the triangle inequality and the fact that projectors have norm 1, we have
\[
\|\hat{P}\| = \|2\Pi_2 + 3\Pi_3\| \le 2\|\Pi_2\| + 3\|\Pi_3\| = 2 + 3 = 5.
\]
For the lower bound, suppose there exists a nonzero vector $x$ such that $\Pi_3 x = x$. Then
\[
\hat{P} x = 2 \Pi_2 x + 3 \Pi_3 x = 2 \Pi_2 x + 3x,
\]
and hence
\[
\|\hat{P}\| \ge \frac{\|\hat{P}x\|}{\|x\|} \ge \frac{\|3x\| - 2\|\Pi_2 x\|}{\|x\|} \ge 3 - 2 \frac{\|\Pi_2 x\|}{\|x\|}.
\]
In particular, choosing $x$ in the range of $\Pi_3$ and orthogonal to the range of $\Pi_2$ (if such a vector exists) gives $\Pi_2 x=0$ and hence $\|\hat{P}\|\ge 3$.
\end{proof}

\subsection{Bounds for prime-attention dynamics}

Recall the prime-attended effective evolution operator
\[
\Xi(t) = \sum_{p\in P} w_p(t) U_p,
\]
where $P$ is a finite prime set, each $U_p$ is unitary, and $w_p(t)\ge 0$, $\sum_p w_p(t)=1$.

\begin{lemma}[Norm bound for $\Xi(t)$]
\label{lem:XiNormBound}
For every $t$,
\[
\|\Xi(t)\| \le 1.
\]
\end{lemma}

\begin{proof}
We use the triangle inequality and Lemma~\ref{lem:UnitaryBound}:
\[
\|\Xi(t)\| = \left\|\sum_{p\in P} w_p(t) U_p\right\|
\le \sum_{p\in P} w_p(t) \|U_p\|
= \sum_{p\in P} w_p(t) = 1.
\]
Thus $\|\Xi(t)\|\le 1$.
\end{proof}

\begin{corollary}[Non-expansiveness of the step map]
Let $T_t:\mathcal{H}\to\mathcal{H}$ be the nonlinear map
\[
T_t(x) = \frac{\Xi(t) x}{\|\Xi(t)x\|},\quad x\neq 0,
\]
and $T_t(0)=0$. Then for any $x$,
\[
\|\Xi(t)x\| \le \|x\|.
\]
\end{corollary}

\begin{proof}
Immediate from $\|\Xi(t)\|\le 1$, since
\[
\|\Xi(t)x\| \le \|\Xi(t)\|\,\|x\| \le \|x\|.
\]
\end{proof}

This does not imply that $T_t$ itself is non-expansive (due to the renormalization), but it shows that the raw step $\Xi(t)$ never increases the norm of a state vector.

\subsection{Prime-attention layer: Lipschitz constants and bounds}

We now derive basic bounds for the prime-attention neural layer
\[
h(x) = \sum_{p\in P} w_p(x)\,\tilde{z}_p
\]
with $\tilde{z}_p = \sigma(W_p x + b_p)$, where $\sigma$ is a 1-Lipschitz activation (e.g.\ ReLU) on $\mathbb{R}^{d_{\mathrm{out}}}$, and $w_p(x)$ are attention weights built from quadratic overlaps.

\subsubsection{Setup}

Fix $P=\{p_1,\dots,p_m\}$, and for each prime $p$:
\begin{itemize}[nosep]
  \item $W_p:\mathbb{R}^{d_{\mathrm{in}}}\to \mathbb{R}^{d_{\mathrm{out}}}$ is linear with operator norm $\|W_p\|$.
  \item $U_p\in\mathbb{R}^{d_{\mathrm{out}}\times k_p}$ has orthonormal columns; $\Pi_p = U_p U_p^\top$ is an orthogonal projector.
\end{itemize}

For input $x\in\mathbb{R}^{d_{\mathrm{in}}}$,
\begin{align*}
z_p(x) &= W_p x + b_p,\\
\tilde{z}_p(x) &= \sigma(z_p(x)),\\
c_p(x) &= \Pi_p \tilde{z}_p(x),\\
a_p(x) &= \|c_p(x)\|_2^2,\\
w_p(x) &= \frac{\exp(\beta a_p(x))}{\sum_{q\in P} \exp(\beta a_q(x))},\\
h(x) &= \sum_{p\in P} w_p(x)\,\tilde{z}_p(x),
\end{align*}
with $\beta>0$ a fixed inverse temperature parameter.

\subsubsection{Operator norms of basic components}

\begin{lemma}[Bounds on branch maps]
\label{lem:BranchLipschitz}
Assume $\sigma$ is 1-Lipschitz (e.g.\ ReLU, $\tanh$). Then, for each $p$,
\[
\|\tilde{z}_p(x) - \tilde{z}_p(y)\|_2 \le \|W_p\|\,\|x-y\|_2,\qquad \forall x,y\in\mathbb{R}^{d_{\mathrm{in}}}.
\]
\end{lemma}

\begin{proof}
We have
\[
\tilde{z}_p(x) - \tilde{z}_p(y)
= \sigma\big(W_p x + b_p\big) - \sigma\big(W_p y + b_p\big).
\]
Using 1-Lipschitzness of $\sigma$,
\[
\|\tilde{z}_p(x) - \tilde{z}_p(y)\|_2
\le \|W_p x - W_p y\|_2
\le \|W_p\|\,\|x-y\|_2.
\]
\end{proof}

\begin{lemma}[Bounds on fibers and overlaps]
\label{lem:FiberOverlapBounds}
For each $p$ and all $x,y$,
\begin{align*}
\|c_p(x) - c_p(y)\|_2 &\le \|W_p\|\,\|x-y\|_2,\\
|a_p(x) - a_p(y)| &\le 2 M_p \|W_p\| \,\|x-y\|_2,
\end{align*}
where $M_p := \sup_{x} \|\tilde{z}_p(x)\|_2$ (assumed finite on the domain of interest).
\end{lemma}

\begin{proof}
We first note that $\Pi_p$ is a projector with $\|\Pi_p\|=1$, since it is an orthogonal projector. Thus
\[
\|c_p(x) - c_p(y)\|_2 = \|\Pi_p \tilde{z}_p(x) - \Pi_p \tilde{z}_p(y)\|_2
\le \|\Pi_p\|\,\|\tilde{z}_p(x) - \tilde{z}_p(y)\|_2
\le \|W_p\|\|x-y\|_2,
\]
using Lemma~\ref{lem:BranchLipschitz}. For the overlaps, we write
\[
a_p(x) = \|c_p(x)\|_2^2 = \langle c_p(x), c_p(x)\rangle.
\]
Then
\begin{align*}
|a_p(x) - a_p(y)|
&= |\|c_p(x)\|_2^2 - \|c_p(y)\|_2^2|\\
&= |\langle c_p(x), c_p(x)\rangle - \langle c_p(y), c_p(y)\rangle|\\
&= |\langle c_p(x)+c_p(y), c_p(x)-c_p(y)\rangle|\\
&\le \|c_p(x)+c_p(y)\|_2 \,\|c_p(x)-c_p(y)\|_2.
\end{align*}
Assuming $\|\tilde{z}_p(\cdot)\|_2\le M_p$ uniformly on the domain we care about, and noting $\|c_p(\cdot)\|\le \|\tilde{z}_p(\cdot)\|$, we have
\[
\|c_p(x)+c_p(y)\|\le \|c_p(x)\| + \|c_p(y)\| \le 2M_p.
\]
Using the bound on $\|c_p(x)-c_p(y)\|$ gives
\[
|a_p(x)-a_p(y)| \le 2 M_p\,\|W_p\|\,\|x-y\|_2.
\]
\end{proof}

\subsubsection{Lipschitz control of attention weights}

The attention weights arise from a softmax applied to the vector $a(x) = (a_p(x))_{p\in P}$. We can bound their sensitivity using standard properties of the softmax.

\begin{lemma}[Softmax Lipschitzness]
\label{lem:SoftmaxLipschitz}
Let $s:\mathbb{R}^m\to\mathbb{R}^m$ be the softmax map
\[
s_j(z) = \frac{\exp(z_j)}{\sum_{k=1}^m \exp(z_k)}.
\]
Then, with respect to the Euclidean norm, there exists a constant $L_s\le 1$ such that
\[
\|s(z)-s(z')\|_2 \le L_s \|z-z'\|_2,\qquad \forall z,z'\in\mathbb{R}^m.
\]
In particular, since each partial derivative of $s$ is bounded by $1/4$ in magnitude, one may take $L_s\le 1$.
\end{lemma}

\begin{proof}
One can bound the Jacobian of $s$ and then bound its operator norm uniformly on $\mathbb{R}^m$. Each entry is
\[
\frac{\partial s_j}{\partial z_k}(z) = s_j(z)(\delta_{jk} - s_k(z)).
\]
It is well-known that the softmax is 1-Lipschitz under the $\ell_2$ norm up to a constant, and that in practice the Jacobian operator norm is bounded by 1; any finite bound suffices for our purposes. A more careful argument can be constructed by estimating the spectral radius of the Jacobian as in standard analysis of softmax.
\end{proof}

Combining Lemmas~\ref{lem:FiberOverlapBounds} and~\ref{lem:SoftmaxLipschitz}, we get:

\begin{proposition}[Lipschitz bound for attention weights]
\label{prop:WeightLipschitz}
Fix $\beta>0$. Let $w(x) = (w_p(x))_{p\in P}$ be prime attention weights as above. Assume $M_p<\infty$ for all $p$. Then there exists a constant
\[
L_w \le \beta L_s \sqrt{|P|}\,\max_{p\in P}(2 M_p \|W_p\|)
\]
such that
\[
\|w(x)-w(y)\|_2 \le L_w \|x-y\|_2,\qquad \forall x,y.
\]
\end{proposition}

\begin{proof}
Let $a(x)=(a_p(x))_{p\in P}$. By Lemma~\ref{lem:FiberOverlapBounds},
\[
\|a(x)-a(y)\|_2 \le \left(\sum_{p\in P} (2M_p\|W_p\|)^2\right)^{1/2} \|x-y\|_2
\le \sqrt{|P|}\,\max_p(2M_p\|W_p\|)\,\|x-y\|_2.
\]
The attention weights are $w(x) = s(\beta a(x))$, where $s$ is softmax. Using the Lipschitz bound for $s$,
\[
\|w(x)-w(y)\|_2
= \|s(\beta a(x)) - s(\beta a(y))\|_2
\le L_s \|\beta a(x) - \beta a(y)\|_2
= \beta L_s \|a(x)-a(y)\|_2,
\]
so
\[
\|w(x)-w(y)\|_2 \le \beta L_s \sqrt{|P|}\,\max_p(2M_p\|W_p\|)\,\|x-y\|_2.
\]
\end{proof}

\subsubsection{Lipschitz bound for the prime-attention layer}

We now bound the Lipschitz constant of $h(x) = \sum_p w_p(x)\tilde{z}_p(x)$.

\begin{theorem}[Global Lipschitz bound for $h$]
\label{thm:LayerLipschitz}
Under the assumptions above (1-Lipschitz activation, finite $M_p$), the prime-attention layer $h$ is globally Lipschitz with
\[
\|h(x)-h(y)\|_2 \le L_h \,\|x-y\|_2,
\]
where one may take, for example,
\[
L_h = \sum_{p\in P} \|W_p\| + \sqrt{|P|}\,\max_{p\in P}\|W_p\| \cdot L_w.
\]
In particular, $L_h$ is finite whenever all $\|W_p\|$ and $M_p$ are finite.
\end{theorem}

\begin{proof}
We write
\begin{align*}
h(x) - h(y)
&= \sum_{p\in P} w_p(x)\tilde{z}_p(x) - \sum_{p\in P} w_p(y)\tilde{z}_p(y)\\
&= \sum_{p\in P} \big( w_p(x)\tilde{z}_p(x) - w_p(x)\tilde{z}_p(y)\big)
 + \sum_{p\in P} \big( w_p(x)\tilde{z}_p(y) - w_p(y)\tilde{z}_p(y)\big)\\
&=: A + B.
\end{align*}

For the term $A$,
\[
\|A\|_2 \le \sum_{p\in P} |w_p(x)|\,\|\tilde{z}_p(x)-\tilde{z}_p(y)\|_2
\le \sum_{p\in P} \|\tilde{z}_p(x)-\tilde{z}_p(y)\|_2,
\]
since $w_p(x)\in[0,1]$. Using Lemma~\ref{lem:BranchLipschitz},
\[
\|A\|_2 \le \sum_{p\in P} \|W_p\|\,\|x-y\|_2.
\]

For $B$,
\[
\|B\|_2 \le \sum_{p\in P} |w_p(x)-w_p(y)|\,\|\tilde{z}_p(y)\|_2
\le \left(\sum_{p\in P} |w_p(x)-w_p(y)|^2\right)^{1/2}
     \left(\sum_{p\in P} \|\tilde{z}_p(y)\|_2^2\right)^{1/2}
\]
by Cauchy--Schwarz. The second factor is bounded by $\sqrt{|P|}\,\max_p M_p$. By Proposition~\ref{prop:WeightLipschitz}, the first factor is bounded by $L_w\|x-y\|_2$. Thus
\[
\|B\|_2 \le L_w \sqrt{|P|}\,\max_p M_p \,\|x-y\|_2.
\]

Combining,
\[
\|h(x)-h(y)\|_2 \le \left(\sum_{p\in P} \|W_p\| + L_w\sqrt{|P|}\,\max_p M_p\right)\|x-y\|_2.
\]
This defines a finite Lipschitz constant $L_h$.
\end{proof}

\subsection{A simple phase-transition heuristic for attention}

While a full rigorous treatment of phase transitions in $w_p$ as $\beta$ varies requires additional assumptions on the distribution of $a_p(x)$, we can at least establish the limiting behavior.

\begin{proposition}[Zero-temperature limit of attention]
\label{prop:ZeroTemp}
Fix $x$ and let $a_p(x)$ be finite for all $p\in P$. Let $\beta\to\infty$. Then
\[
w_p(x) = \frac{\exp(\beta a_p(x))}{\sum_{q\in P} \exp(\beta a_q(x))}
\]
converges to a one-hot distribution supported on the set of maximizers of $a_p(x)$:
\[
\lim_{\beta\to\infty} w_p(x) =
\begin{cases}
\frac{1}{|\arg\max_q a_q(x)|}, &\text{if } a_p(x) = \max_q a_q(x),\\
0, &\text{otherwise.}
\end{cases}
\]
\end{proposition}

\begin{proof}
Let $m = \max_q a_q(x)$ and write $a_p(x) = m - \delta_p$, where $\delta_p\ge 0$ and $\delta_p=0$ iff $p$ is a maximizer. Then
\[
w_p(x) = \frac{\exp(\beta (m - \delta_p))}{\sum_q \exp(\beta (m - \delta_q))}
= \frac{\exp(-\beta \delta_p)}{\sum_q \exp(-\beta \delta_q)}.
\]
For any $p$ with $\delta_p>0$, $\exp(-\beta \delta_p)\to 0$ as $\beta\to\infty$. For any maximizer $p$ with $\delta_p=0$, $\exp(-\beta \delta_p)=1$ for all $\beta$. Let $S = \{p:\delta_p=0\}$ be the set of maximizers. Then as $\beta\to\infty$,
\[
\sum_q \exp(-\beta\delta_q) \to |S|,
\]
and thus $w_p(x)\to 1/|S|$ for $p\in S$ and $w_p(x)\to 0$ for $p\notin S$.
\end{proof}

This shows that as $\beta$ grows, the system approaches a ``hard'' prime routing regime where attention effectively selects argmax primes. The regularizer $\mathcal{L}_{\mathrm{sparse}}$ encourages this regime even at finite $\beta$.

\bigskip

These basic estimates provide initial analytical control over the objects appearing in the MPD models and justify, at least in finite settings, that the constructions are well-behaved and amenable to further functional-analytic and dynamical analysis.

\nocite{*}
\bibliographystyle{plainnat}
\bibliography{references}
\end{document}